\chapter{Theoretical Framework and Operationalisation}
\label{chap:theory}

\section{Conceptual framework}
\label{sec:theory}

\subsection{Purpose and analytical levels}

This thesis explains how factual content changes while it moves through a synthetic social network.
The explanation distinguishes three levels: persona composition determines which agents transform a message, topology determines potential exposure, and propagation links agent transformations and exposure over time.
The resulting texts are then scored with an explicit measurement instrument.

This sequence separates mechanism from measurement.
The simulated mechanism produces texts, actions, and exposure paths.
The auditor converts those texts into misinformation scores.
The scores are observations of the simulation.
They are not the mechanism itself.

The framework also separates misinformation from a firestorm.
Misinformation concerns departure from factual source claims.
An online firestorm concerns a rapid and large discharge of negative word-of-mouth \parencite{pfeffer2014firestorms}.
The two phenomena may coincide.
They need not coincide.
Accurate criticism can form a firestorm.
Severe misinformation can also diffuse slowly and quietly.
Therefore, Misinformation Index (MI) and Misinformation Propagation Rate (MPR) do not by themselves identify a firestorm.

Three kinds of statement follow from this distinction.
\emph{Theory} states a proposed relation among composition, topology, exposure, and content change.
\emph{Operationalisation} states how the relation is represented in the simulation.
\emph{Evidence} states what an observed source or experiment can support.
The distinction is maintained throughout this chapter.

\subsection{The integrated generative argument}
\label{subsec:integrated-causal-argument}

Let \(G=(V,E,W)\) be a directed network.
The node set \(V\) contains synthetic agents.
The edge set \(E\) defines possible transmissions.
The weights \(W\) represent relational conditions such as trust.
Each node \(i\) receives a persona \(p_i\).
The complete assignment is denoted by \(\boldsymbol{p}=(p_1,\ldots,p_{|V|})\).

Let \(X_e\) be the text carried by propagation event \(e\).
When node \(i\) processes the event at time \(t\), it applies a persona-conditioned transformation:
\begin{equation}
X_{e'} = T_{p_i}\!\left(X_e,\,H_i(t),\,R_i(t)\right).
\label{eq:persona-transform}
\end{equation}
Here, \(H_i(t)\) is the information available to the agent.
\(R_i(t)\) denotes local relational conditions.
The output may be forwarded, reinterpreted, or dropped.
Thus, propagation is a sequence of selective text transformations.
It is not the movement of a fixed information token.

The theoretical object is the joint process
\begin{equation}
\mathcal{Y}
=
F\!\left(
\boldsymbol{p},
G,
S,
A,
\mathcal{T},
\varepsilon
\right),
\label{eq:joint-process}
\end{equation}
where \(S\) is the seed message.
\(A\) is the action rule.
\(\mathcal{T}\) is the exposure schedule.
\(\varepsilon\) collects stochastic and model-specific variation.
\(\mathcal{Y}\) contains generated texts, transmissions, and cascade structure.
An auditor \(J\) then maps the generated texts to observed scores:
\begin{equation}
\widehat{\mathcal{M}} = J(\mathcal{Y};Q,G_Q,\theta_J).
\label{eq:measurement-map}
\end{equation}
The question set \(Q\), gold answers \(G_Q\), and auditor specification \(\theta_J\) are part of the measurement model.
They are not neutral bookkeeping choices.

The framework therefore has a clear order.
Composition and topology shape exposure.
Exposure selects which persona transformations occur.
Repeated transformations produce factual retention or drift.
The auditor then measures that drift.
Pfeffer's firestorm theory interprets the broader communication conditions.
Debnath's belief profiles ground the identities and issue domain.

\subsection{Persona composition as a transformation mechanism}
\label{subsec:persona-composition}

\paragraph{Theory.}
A persona is a structured prior over attention, interpretation, and expression.
It makes some source claims salient.
It makes other claims easy to omit or recode.
Repeated use of the same prior can compound a stable distortion.
Mixed priors can interrupt that process.
They can also introduce new distortions.
Heterogeneity is therefore not intrinsically corrective.

Composition has at least three dimensions.
The first is the share of each persona family.
The second is the diversity of those families.
The third is their placement in the graph.
Two societies can have the same shares but different outcomes.
One may place a conspiracy persona at a hub.
The other may place the same persona at the periphery.
The label ``heterogeneous'' does not capture this difference.

Results from the prior auditor--node preprint illustrate this dependence on composition and order.
In homogeneous 30-node branches, identity- and ideology-laden personas often accelerated drift.
Expert and neutral personas often stabilized factual content \parencite{maurya2025simulating}.
Yet the controlled-random heterogeneous branches reached the high-loss band
(\(\mathrm{MI}>3\)) in about \(85\%\) of branch--domain pairs.
The prior study therefore does not establish a general ``diversity buffers misinformation'' law.
It shows that persona effects depend on order, content, and repeated interaction.

\paragraph{Operationalisation.}
The thesis assigns persona identifiers and prompts to nodes.
Homogeneous conditions repeat a focal identity.
Heterogeneous conditions combine several identities.
The relevant treatment is the complete assignment \(\boldsymbol{p}\), not the H/He label alone.
Composition should therefore be reported as persona-family shares and node positions.
Results should also be conditioned on persona family.
A lower network mean may otherwise reflect fewer conspiracy seats rather than corrective interaction.

\paragraph{Evidence.}
The prior arXiv preprint provides simulation evidence for persona-conditioned drift on linear branches.
It does not provide human behavioural validation.
It also does not provide evidence about complex-network placement.
The present experiments can provide evidence about synthetic agents under fixed prompts and models.
They cannot establish that real users with similar labels would behave identically.

\subsection{Topology as an exposure mechanism}
\label{subsec:topology-exposure}

\paragraph{Theory.}
Topology affects factual content indirectly by determining which agents can receive and transform a message.
Degree controls how many transmissions can follow an event.
Path length controls how many sequential transformations can occur.
Clustering creates repeated exposure from connected neighbours.
Community structure concentrates exposure within groups.
Bridges permit a frame to cross between groups.
Hubs place selected agents in high-leverage positions.

Topology and composition are therefore non-separable.
A clustered graph does not create an identity echo unless identities align with its clusters.
A homogeneous graph produces perfect identity agreement by construction.
That agreement is not evidence of an emergent echo chamber.
Likewise, a scale-free graph is primarily a hub contrast.
It is not automatically a clustering contrast.

\paragraph{Operationalisation.}
The topology generator defines \(E\) and initial edge conditions.
Persona assignment then defines the relation between social position and identity.
The design includes chains, rings, random graphs, small-world graphs, scale-free graphs, echo chambers, polarised graphs, and hierarchical graphs.
Each topology represents a distinct exposure opportunity structure.
No generator is treated as a complete model of a platform.

The linear chain is the direct link to the prior arXiv preprint.
The small-world graph is the clearest named operationalisation of local clustering \parencite{watts1998collective}.
The scale-free graph represents unequal connectivity and hub exposure \parencite{barabasi1999emergence}.
The random graph provides a stochastic baseline \parencite{erdos1960evolution}.
The echo and polarised graphs encode pre-existing segregation.
They do not demonstrate that segregation emerged during the run.
The hierarchical graph represents asymmetric transmission.
It does not operationalise legacy-media feedback.

\paragraph{Evidence.}
Differences between topology cells are evidence about the implemented generators.
They are not evidence about all networks in the same named class.
One archived realisation per cell does not identify a topology-wide effect.
Dead cascades also carry information about reachability.
They must not be recoded as zero misinformation.

\subsection{Propagation as a coupled selection process}
\label{subsec:propagation}

\paragraph{Theory.}
Propagation determines which potential transformations become realized.
At each event, an agent may forward, reinterpret, or drop a message.
Forwarding preserves more of the received text.
Reinterpretation invokes another persona-conditioned rewrite.
Dropping ends that path.
The observed cascade is therefore selected by both network access and agent action.

This creates two distinct margins.
The extensive margin concerns whether content reaches another node.
The intensive margin concerns how much the content changes when it is processed.
A topology can increase reach without increasing factual drift per rewrite.
A persona can increase drift while reducing onward transmission.
These mechanisms should not be collapsed into one mean.

Repeated exposure adds a further mechanism.
Several neighbours may deliver related frames to the same node.
The node can then receive social confirmation as well as content.
Pfeffer describes this local echo as a consequence of network clusters \parencite{pfeffer2014firestorms}.
In the simulation, however, repeated receipt is an implemented exposure pattern.
It is not direct evidence of human persuasion.

\paragraph{Operationalisation.}
The simulator records event time, hop count, sender, receiver, action, and rewritten text.
Ticks and hops answer different questions.
A tick is a simulation update.
A hop is a position along a transmission path.
Neither is a unit of wall-clock time.
The eight-tick horizon is a compressed experimental horizon.
It is not an estimate of a Twitter half-life.
It is also not Debnath's eight-word skip-gram window.

The action set is ternary.
It includes forwarding, reinterpretation, and dropping.
This differs from Pfeffer's binary-choice factor.
No result from the ternary mechanism can confirm a binary-choice proposition without a dedicated ablation.

\paragraph{Evidence.}
Propagation logs support claims about simulated reach, depth, breadth, and rewriting.
They do not support claims about platform speed in minutes or hours.
They also do not identify causal persuasion by repeated exposure.
Such claims would require a matched manipulation of exposure while other conditions remain fixed.

\subsection{MI and MPR as a measurement model}
\label{subsec:mi-mpr}

\paragraph{Theory.}
Factual drift is latent until a measurement rule is applied.
The auditor--node framework makes that rule explicit \parencite{maurya2025simulating}.
The source article defines a set of factual questions.
The auditor evaluates whether each answer remains recoverable from a rewrite.
MI summarizes deviation at one scored event.
MPR summarizes MI over a specified collection of events.

In the prior arXiv preprint, the answer vector has ten binary items.
MI is the Hamming count between the reference and audited vectors.
Branch MPR is the mean MI across a 30-node branch.
Despite its name, MPR is not a temporal derivative.
It is a branch-level average severity measure.
This matters when the concept is transferred to a graph.

\paragraph{Operationalisation.}
For discrete scoring with \(m\) questions, let \(z_{ej}=1\) when answer \(j\) is not preserved at event \(e\).
Then
\begin{equation}
\operatorname{MI}^{(d)}_e
=
\sum_{j=1}^{m} z_{ej}.
\label{eq:discrete-mi}
\end{equation}
The present climate instrument uses \(m=5\).
Its numerical range therefore differs from the prior arXiv preprint's ten-item instrument.

On a graph, there is no unique branch average.
The aggregation set must be named.
A node summary can be written as
\begin{equation}
\operatorname{MPR}_i
=
\frac{1}{|E_i^{\mathrm{scored}}|}
\sum_{e\in E_i^{\mathrm{scored}}}
\operatorname{MI}_e.
\label{eq:node-mpr}
\end{equation}
A network-time summary can instead average scored events at tick \(t\):
\begin{equation}
\overline{\operatorname{MI}}(t)
=
\frac{1}{|E_t^{\mathrm{scored}}|}
\sum_{e\in E_t^{\mathrm{scored}}}
\operatorname{MI}_e.
\label{eq:network-mi}
\end{equation}
These quantities answer different questions.
Neither should be labelled simply ``MPR'' without its aggregation rule.

The high-loss threshold is also operational.
The prior arXiv preprint labels averages at or below \(1\) as factual error.
It labels values above \(1\) and at or below \(3\) as lies.
It labels values above \(3\) as a high-loss or ``propaganda'' band in that
ten-item instrument \parencite{maurya2025simulating}. The present five-item
reuse of \(>3\) is an arbitrary sensitivity cut, not a calibrated propaganda
category.
The thesis reuses these cuts.
They are design thresholds.
They are not validated psychometric boundaries.

\paragraph{Evidence.}
MI supports a claim about auditor-detected recoverability of selected source claims.
It does not directly measure belief, intent, toxicity, affect, or public harm.
MPR supports a claim about the average of those event scores over a declared set.
It does not show that misinformation propagated faster.
The prior arXiv preprint assumes a perfect auditor in its limitations.
The present thesis must instead treat auditor error as a validity threat.

\subsection{Instrument dependence}
\label{subsec:instrument-dependence}

\paragraph{Theory.}
Observed misinformation is jointly produced by the generated text and the instrument.
Formally,
\begin{equation}
\widehat{M}_{e,r}
=
J_r(X_e,Q,G_Q),
\label{eq:instrument-index}
\end{equation}
where \(r\) indexes an auditor mode.
Changing \(r\) can change levels, rankings, and threshold crossings.
The latent text \(X_e\) may remain unchanged.
Therefore, a score difference across instruments is not automatically a content difference.

\paragraph{Operationalisation.}
The continuous campaign and the dual-discrete campaign are separate measurement regimes.
The continuous auditor returns graded claim-level assessments.
The discrete headline counts claim failures.
The continuous sidecar in a dual run is a diagnostic attached to that regime.
It is not a third ground truth.
Scores from different regimes must not be averaged.
Threshold crossings must be computed within one regime.

Instrument dependence has four sources.
First, question selection determines which facts can be lost.
Second, response coding determines whether partial preservation receives credit.
Third, the auditor model introduces judgement variance and bias.
Fourth, aggregation determines which nodes and events receive weight.
Topology can affect the fourth source by changing event counts.
An event-weighted mean can therefore partly reflect cascade selection.

\paragraph{Evidence.}
Agreement across instruments strengthens robustness.
Disagreement identifies a measurement sensitivity.
It does not reveal which instrument is true.
External validation requires human coding or another justified reference.
LLM-as-judge studies motivate automated assessment, but they do not validate this climate-specific instrument by transfer alone \parencite{zheng2023judge}.

\subsection{From branches to complex networks}
\label{subsec:complex-network-generalization}

\paragraph{Theory.}
The prior arXiv preprint studies sequential transmission chains.
Each node has one predecessor in the realized branch.
Depth and exposure order are therefore aligned.
A complex network breaks this alignment.
A node can receive several messages.
A message can follow several paths.
Communities, bridges, hubs, and cycles can alter both reach and rewrite order.

The extension is conceptual, not merely technical.
In a chain, composition is an ordered list.
In a graph, composition becomes a joint distribution over identities and positions.
In a chain, MPR has one natural branch denominator.
In a graph, several valid denominators exist.
In a chain, a late node implies many prior rewrites.
In a graph, a late tick need not imply a long path.

The generalization therefore preserves the event-level MI concept.
It does not preserve every branch-level interpretation.
Graph analyses must report path, node, event, and time aggregation separately.

\paragraph{Operationalisation.}
The chain serves as a bridge to prior work.
The topology grid then changes the adjacency structure while retaining the same basic rewrite and audit process.
Persona composition is crossed with topology where the design permits.
This crossing tests whether identity effects depend on exposure structure.
It does not constitute a literal replication of the 30-hop study.
The climate articles, persona inventory, model, horizon, and graph process differ.

Comparisons also require common support.
Reach, event count, and dead-cell status should accompany MI summaries.
A live-only mean describes surviving cascades.
It does not estimate the mean over cascades that failed to propagate.
Similarly, a 63-node imported graph should not be pooled with an eight-node generated graph.
The opportunity structures differ too sharply.

\paragraph{Evidence.}
The complex-network experiments can show whether a prior pattern survives a domain and structure shift.
They cannot reproduce the prior arXiv preprint merely by reusing MI and MPR.
Evidence of generalization requires directional stability under the changed object.
Evidence of failure is also informative.
It places a boundary on the chain result.

\subsection{Debnath belief profiles as empirical grounding}
\label{subsec:debnath-profiles}

\paragraph{Theory.}
Debnath et al.\ describe geoengineering discourse as an identity-structured information environment \parencite{debnath2023conspiracy}.
Chemtrails conspiracy content dominates the observed corpus.
Smaller streams concern climate action and environmental risk.
These streams provide a domain-specific basis for synthetic belief profiles.
They also motivate identity alignment between a message and a persona.

The belief profiles are mechanisms of selective interpretation.
A conspiracy profile is expected to assimilate ambiguous geoengineering information into a chemtrails frame.
A climate-action profile may foreground governance and justice.
An environmental profile may foreground ecological risk.
Added scientist, journalist, and moderation roles supply institutional accuracy norms.
These added roles support the experimental contrast.
They are not empirical Debnath clusters.

\paragraph{Operationalisation.}
The persona library translates discourse types into prompts, tags, tones, and expertise fields.
This is a theory-faithful reduction.
It is not a reconstruction of individual Twitter users.
The full tweet corpus was not hydrated.
HDBSCAN profiles were not re-estimated.
The imported 63-node structure is a hashtag co-occurrence graph.
It is not a retweet cascade.

Debnath's skip-gram context window concerns words inside tweets.
It does not define simulation hops.
Debnath also reports no empirical MI or MPR.
Twitter toxicity, semantic proximity, and simulated factual drift remain different constructs.

\paragraph{Evidence.}
Debnath provides observational evidence about the issue domain, discourse streams, and platform corpus.
It supports the content validity of the persona families.
It does not validate their exact prompts.
It also does not validate simulated propagation rates.
Claims about the synthetic profiles must therefore remain claims about grounded experimental agents.

\subsection{Pfeffer firestorms as the macro-theoretical frame}
\label{subsec:pfeffer-firestorms}

\paragraph{Theory.}
Pfeffer, Zorbach, and Carley define online firestorms as sudden, large discharges of negative word-of-mouth \parencite{pfeffer2014firestorms}.
Their Outlook identifies seven interrelated factors:
\begin{enumerate}
    \item speed and volume of communication;
    \item binary choices;
    \item network clusters;
    \item unrestrained information flow;
    \item lack of diversity;
    \item cross-media dynamics; and
    \item network-triggered decision processes.
\end{enumerate}
This list is the theoretical source.
It is not a list of simulator parameters.

The framework maps these factors onto different analytical levels.
Speed and volume concern propagation intensity.
Binary choices concern the action architecture.
Clusters and unrestrained flow concern topology.
Lack of diversity concerns composition and homophily.
Network-triggered decisions concern repeated social exposure.
Cross-media dynamics concerns an external feedback layer.
The affective character of firestorms belongs to the definition.
It is not an eighth numbered factor.

\paragraph{Operationalisation.}
The thesis uses a computational remapping.
Identity composition operationalises part of lack of diversity.
Graph generators operationalise selected cluster and flow conditions.
Seed fan-out can operationalise shock versus drip.
Ticks and activity schedules can operationalise simulated opportunity rates.
These mappings are partial.

Two factors remain especially constrained.
The action architecture is ternary rather than binary.
The model has no legacy-media broadcast layer.
Cross-media dynamics is therefore held outside the simulation.
The hierarchical topology cannot substitute for that missing layer.

Valence, surprise, identity, clustering, echo, and temporal acceleration are thesis-level knobs.
They are not Pfeffer's original factor names.
Their value lies in making parts of the theory manipulable.
Their limitation lies in incomplete construct coverage.

\paragraph{Evidence.}
Pfeffer's article provides qualitative and case-based theoretical evidence.
It does not provide an LLM model, MI, MPR, or a numerical tipping threshold.
An MI value above \(3\) is therefore not a Pfeffer firestorm criterion.
Simulation evidence can test patterns implied by the remapping.
It cannot confirm all seven factors when some factors are held or absent.

\subsection{A two-axis outcome model}
\label{subsec:two-axis-model}

The integrated framework uses two outcome axes.
The first axis is communication escalation.
It includes reach, volume, breadth, depth, and simulated speed.
The second axis is content degradation.
It includes event MI and declared MPR aggregates.
Firestorm interpretation also requires negative or affective content.

This yields four conceptual cases (\Cref{tab:two-axis-outcomes}).
\begin{table}[t]
\centering
\caption{Communication escalation and factual degradation are distinct outcomes.}
\label{tab:two-axis-outcomes}
\begin{tabular}{p{0.19\textwidth}p{0.34\textwidth}p{0.34\textwidth}}
\toprule
 & Low factual degradation & High factual degradation \\
\midrule
Low communication escalation
& Limited, largely faithful diffusion
& Localized misinformation without a firestorm \\
High communication escalation
& Viral but factually stable communication; it may still be an affective firestorm
& A misinformation-amplifying firestorm candidate, if negative affect is also present \\
\bottomrule
\end{tabular}
\end{table}

The final cell is the thesis's central concern.
It is still a candidate category.
The simulation must show escalation and degradation separately.
It must also justify the affective interpretation.
No single mean MI can establish all three conditions.

\subsection{Derived propositions}
\label{subsec:derived-propositions}

The framework yields five propositions.
P1 and P5 are examined directly in the present campaign.
P2--P4 identify mechanisms that the current bundled design can describe only
indirectly and that require matched ablations for a direct test.
They do not report findings.

\paragraph{P1: composition--alignment proposition.}
Factual drift should increase when a persona's interpretive prior aligns with a misinformation frame.
The effect should be strongest under repeated use of that persona.
However, heterogeneous composition may either interrupt or compound drift.
Its direction depends on order and role composition.

\paragraph{P2: position--composition proposition.}
The effect of a persona depends on network position.
Placing an aligned persona at a hub should alter more downstream opportunities than placing it at the periphery.
Composition shares alone are therefore insufficient.

\paragraph{P3: cluster--echo proposition.}
Clustering should increase repeated local exposure.
Identity echo should increase only when persona placement aligns with the clustered structure.
Structural clustering and identity homophily are distinct conditions.

\paragraph{P4: reach--severity proposition.}
Topology may change cascade reach without changing drift per scored rewrite.
Persona transformations may change drift without increasing reach.
Communication escalation and content degradation should therefore be analysed separately.

\paragraph{P5: instrument-conditional proposition.}
The magnitude and threshold status of measured drift depend on the auditor regime.
A substantive conclusion is robust only when its direction survives justified measurement alternatives.
Instrument disagreement is a result about measurement.
It is not evidence of two underlying cascades in the same run.

\subsection{Theory, operationalisation, and evidence ledger}
\label{subsec:theory-ledger}

The ledger in \Cref{tab:theory-operation-evidence} separates theoretical
constructs, operational choices, and admissible evidence.

\begin{table}[tbp]
\centering
\caption{Separation of theoretical constructs, operational choices, and admissible evidence.}
\label{tab:theory-operation-evidence}
\scriptsize
\begin{tabular}{p{0.13\textwidth}p{0.22\textwidth}p{0.24\textwidth}p{0.22\textwidth}}
\toprule
Construct & Theory & Operationalisation & Evidence boundary \\
\midrule
Persona composition
& Distribution and placement of interpretive priors shape text transformation.
& Persona prompts; homogeneous and heterogeneous assignments; family shares; node positions.
& Synthetic-agent evidence only. H/He alone does not identify diversity effects. \\

Topology
& Adjacency, clustering, hubs, and bridges shape exposure opportunities.
& Eight graph generators plus an imported custom graph.
& Generator-specific evidence. One archived realisation does not identify a topology class. \\

Propagation
& Selective forwarding, rewriting, and dropping jointly determine realized paths.
& Ternary actions; ticks; hops; trust; inbox and horizon constraints.
& Simulation time is not wall-clock time. Ternary actions do not test binary choice. \\

Factual drift
& Source claims may be omitted, contradicted, or transformed.
& QA-based discrete or continuous auditor.
& Measures selected claim recoverability, not belief, intent, toxicity, or harm. \\

MI
& Event-level observed factual deviation.
& Five-item climate score within a declared auditor mode.
& Scale depends on question count, coding, and auditor. \\

MPR
& Aggregate severity over a declared propagation unit.
& Branch, node, event, or time mean with an explicit denominator.
& Not a physical rate. Graph aggregates are not automatically comparable with chain MPR. \\

Debnath profiles
& Issue-specific identities condition interpretation of geoengineering claims.
& Theory-faithful persona prompts and coarse discourse families.
& Grounded reductions, not recovered user clusters or digital twins. \\

Pfeffer firestorm
& Rapid, high-volume, negative word-of-mouth under seven interrelated conditions.
& Partial mapping to composition, topology, seeding, and schedule.
& MI/MPR alone cannot identify a firestorm. Cross-media remains untested. \\
\bottomrule
\end{tabular}
\end{table}

\subsection{Scope of inference}
\label{subsec:scope-inference}

The framework supports mechanism-oriented simulation claims.
It can show that a specified persona, graph, seed, and auditor produce a specified pattern.
It can compare matched cells within one measurement regime.
It can test whether chain patterns persist on more complex networks.

The framework does not support direct population inference.
The agents are not sampled humans.
The graphs are not representative platform samples.
The auditor is not an error-free observer.
The imported hashtag graph is not a reconstructed diffusion cascade.
Accordingly, the evidence is computational and conditional.

This limitation does not remove the theoretical contribution.
It clarifies it.
The contribution is an explicit bridge across three literatures.
Pfeffer supplies the macro conditions of online firestorms.
Debnath supplies the climate-conspiracy identities and domain.
Maurya et al.\ supply the persona rewrite and QA-auditor instrument.
The thesis combines these elements in a complex-network design and tests which findings persist across network structures and measurement instruments.
